\documentclass[aspectratio=169,10pt]{beamer}
\usetheme{metropolis}
\metroset{block=fill, sectionpage=progressbar}
\setbeamersize{text margin left=0.7cm, text margin right=0.7cm}
\usepackage{fontspec}
\setmainfont{Noto Serif CJK KR}
\setsansfont{Noto Sans CJK KR}
\setmonofont{DejaVu Sans Mono}
\usepackage{amsmath, amssymb}
\usepackage{xcolor}
\usepackage{hyperref}

\definecolor{aprlblue}{HTML}{1F4E79}
\definecolor{aprllight}{HTML}{4A7FB5}
\definecolor{aprlaccent}{HTML}{D9531E}
\definecolor{aprlgray}{HTML}{4A4A4A}
\definecolor{aprlbg}{HTML}{F2F4F7}

\setbeamercolor{frametitle}{bg=aprlblue, fg=white}
\setbeamercolor{title separator}{fg=aprlaccent}
\setbeamercolor{progress bar}{fg=aprlaccent, bg=aprlbg}
\setbeamercolor{alerted text}{fg=aprlaccent}

\newenvironment{script}{%
  \footnotesize
  \setlength{\parskip}{0.5em}
  \setlength{\parindent}{0pt}
}{}

\hypersetup{colorlinks=true, urlcolor=aprllight}

\title{JEPA 강의 --- 발표 대본}
\subtitle{LLM 너머의 로봇 지능을 향해}
\author{Giseop Kim}
\institute{APRL, DGIST}
\date{}

\begin{document}

% ============================================================
{
\setbeamertemplate{footline}{}
\begin{frame}[plain]
\vspace*{0.8cm}
\begin{center}
{\LARGE\textbf{JEPA 강의 --- 발표 대본}}\\[0.3cm]
{\large LLM 너머의 로봇 지능을 향해}\\[0.8cm]
\textbf{Giseop Kim}\\[0.15cm]
{\small APRL, DGIST}
\end{center}

\vspace{0.8cm}
\begin{block}{출처 (Source)}
본 강의 자료는 다음 영상을 기반으로 재구성되었습니다.
\begin{itemize}
  \footnotesize
  \item \textbf{제목}: \textit{Yann LeCun's Billion Dollar Bet}
  \item \textbf{저자}: \href{https://www.youtube.com/@WelchLabs}{Welch Labs}
  \item \textbf{영상 URL}: \href{https://youtu.be/kYkIdXwW2AE?si=asUnIL7EKPRMRwDp}{\texttt{youtu.be/kYkIdXwW2AE}}
  \item \textbf{추가 출처}: LeCun (2022) ``A Path Towards Autonomous Machine Intelligence''
\end{itemize}
\end{block}
\end{frame}
}

% ============================================================
\begin{frame}{이 문서에 대하여}
\begin{itemize}
  \item 이 문서는 \texttt{lec02\_vla\_lecture.tex}(슬라이드 본편)와 \textbf{1:1로 대응}되는 발표 대본입니다.
  \item 한 슬라이드당 \textbf{2--4분 분량}의 구어체 설명 (총 약 60--90분 강의 분량).
  \item 강의 시 \texttt{lec02\_vla\_lecture.pdf}를 띄워놓고 이 문서를 \textbf{진행 노트}로 활용.
  \item 대본은 그대로 읽기보다는 \alert{흐름과 강조점을 잡는 가이드}로 사용.
\end{itemize}

\vspace{0.8em}
\textbf{표기}
\begin{itemize}
  \item 각 슬라이드 대본 시작 부분에 \textbf{[\,슬라이드 번호 / 제목\,]} 표시
  \item \textit{기울임체}: 청중의 반응을 유도하는 질문 또는 강조 톤
  \item [말줄임]: 호흡을 두는 자리
\end{itemize}
\end{frame}

% ============================================================
\section{Part 1. LeCun의 도발과 JEPA}

\begin{frame}{[3] 10억 달러짜리 베팅: 한 마디에서 시작하다}
\begin{script}
자, 지난 주에는 VLA와 $\pi_0$를 통해 LLM 기반 로봇 지능의 최전선을 봤습니다. 오늘은 정반대의 시각에서 출발해 봅시다.

도발적인 인용으로 시작할게요. 얀 르쿤(Yann LeCun)의 말입니다. \textit{``에이전트 시스템을 만들겠다면서, 자기 행동의 결과를 예측할 능력조차 없는 시스템을 떠올린다는 게 도무지 이해가 안 됩니다.''} 실리콘밸리에 적을 좀 만들 만한 발언이죠? 그런데 그는 진심입니다.

르쿤은 AI 분야에서 Yoshua Bengio, Geoffrey Hinton과 함께 ``딥러닝의 세 선구자'' 중 한 명으로, 2018년 튜링상을 함께 받은 사람이에요. 오랫동안 Meta의 수석 AI 과학자로 있으면서 LLM 시대를 같이 만들어왔습니다. 그런 그가 최근 \textbf{10억 달러를 모아서} 전혀 다른 방향의 AI를 추구하고 있어요.

[말줄임] 그 핵심이 \textbf{JEPA}입니다. Joint Embedding Predictive Architecture. 오늘 이게 무엇인지, 왜 르쿤이 이게 LLM보다 낫다고 믿는지를 추적해 봅시다. 그리고 그 과정에서 자기지도학습, 표현 붕괴, 세계 모델이라는 개념들을 하나의 이야기로 엮어볼 거예요.
\end{script}
\end{frame}

% ============================================================
\begin{frame}{[4] JEPA: 아이디어의 핵심}
\begin{script}
먼저 JEPA가 기존 ML과 어떻게 다른지를 구조 수준에서 봅시다.

기존 머신러닝의 기본 공식은 간단합니다. 입력 $X$가 주어지면 출력 $Y$를 예측하도록 모델을 학습시킨다는 거예요. LLM은 입력 텍스트 $X$ 다음에 올 텍스트 $Y$를 예측하고, 이미지 분류기는 이미지 $X$에 해당하는 라벨 $Y$를 예측합니다. 이건 우리가 잘 아는 패러다임이에요.

JEPA는 이 구도를 비틉니다. $X$와 $Y$ 각각을 \textbf{인코더}에 통과시켜 임베딩, 즉 숫자 벡터로 변환한 다음, 세 번째 모델인 \textbf{예측기}가 $X$의 임베딩으로부터 $Y$의 임베딩을 예측하도록 학습합니다.

[말줄임] 이 차이가 처음에는 작아 보이는데, 굉장히 중요한 함의가 있어요. 첫째, 모델이 픽셀이나 토큰을 직접 \emph{생성}하지 않습니다. 추상화된 표현 공간에서 예측하는 거예요. 둘째, 예측하는 대상이 고차원 픽셀 배열이 아니라 상대적으로 저차원의 \textbf{의미 있는 표현}입니다.

왜 이게 중요한가를 이해하려면, 르쿤이 어떤 역사적 경로를 거쳐서 이 결론에 이르렀는지를 봐야 해요. 지금부터 그 여정을 따라가 봅시다.
\end{script}
\end{frame}

% ============================================================
\section{Part 2. CNN에서 자기지도학습까지}

\begin{frame}{[5] LeCun의 역사: 혁명을 두 번 예감한 사람}
\begin{script}
르쿤은 이미 1980년대에 한 번 혁명을 예감한 사람입니다. 당시 AI 학계 대부분은 \textbf{전문가 시스템}에 매달려 있었어요. 사람이 일일이 규칙을 짜 넣는 방식이죠. 그런데 르쿤은 데이터에서 직접 학습하는 \textbf{합성곱 신경망(CNN)}을 개척했습니다.

그로부터 25년이 지나 딥러닝이 주류로 떠오르며 등장한 AlexNet이 르쿤의 1990년대 CNN과 놀라울 만큼 닮아 있었다는 건 유명한 이야기예요. 2012년에 ImageNet에서 이전 방법들을 압도하면서 딥러닝 시대가 열렸죠.

[말줄임] 그런데 르쿤을 포함한 연구자들은 이 성공의 이면에서 한 가지 불편한 사실을 봤습니다. AlexNet의 성공은 \textbf{사람이 라벨링한 대용량 데이터}에 크게 기대고 있었다는 거예요. ImageNet에는 사람이 손수 라벨을 단 120만 장의 이미지가 있었습니다.

이 문제를 느낀 연구자들이 물었어요. \textit{``아이들은 몇 번만 봐도 강아지를 인식하는데, 왜 우리 모델은 수십만 장의 라벨 이미지가 필요한 거지?''} 이게 자기지도학습에 대한 관심으로 이어지는 출발점입니다.
\end{script}
\end{frame}

% ============================================================
\begin{frame}{[6] 라벨 병목: 지도학습의 한계}
\begin{script}
지도학습에 대한 르쿤의 비판을 좀 더 구체적으로 살펴봅시다.

핵심은 \textbf{라벨이 병목}이라는 거예요. AlexNet 같은 지도학습 모델이 잘 되려면 대용량의 라벨 있는 데이터가 필요한데, 이 라벨을 만드는 게 엄청난 비용과 시간을 필요로 합니다. ImageNet 120만 장을 라벨링하는 데 수년이 걸렸어요.

[말줄임] 반면에 아이를 생각해 봅시다. 아이는 ``강아지''라는 단어를 몇 번 들으면서, 실제 강아지를 몇 번 보는 것만으로 강아지를 인식하는 법을 배웁니다. 아무도 강아지 사진에 ``이게 강아지야''라고 라벨을 붙여서 수만 장을 보여주지 않아요.

이 차이가 시사하는 건, 인간의 학습에는 명시적 라벨 없이도 유용한 표현을 얻는 메커니즘이 있다는 겁니다. 바로 \textbf{자기지도학습(self-supervised learning)}이 그 메커니즘을 모방하려는 시도예요. 데이터 자체에서 학습 신호를 뽑아내는 거죠. 다음 토큰 예측에서 ``다음 토큰''이 바로 데이터 자체에서 나오는 라벨인 것처럼요.

르쿤이 2015년쯤부터 들고 다닌 케이크 슬라이드를 봅시다.
\end{script}
\end{frame}

% ============================================================
\begin{frame}{[7] LeCun의 케이크 슬라이드 (2015$\sim$)}
\begin{script}
이 케이크 비유가 유명한 이유는, 당시 학계의 흐름과 \textbf{정면으로 반대되는 주장}이었기 때문입니다.

르쿤은 이렇게 말했어요. \textit{``지능이 케이크라면, 케이크의 본체는 자기지도학습입니다. 위에 얹는 아이싱이 지도학습이고, 그 위에 꽂힌 체리가 강화학습이죠.''}

2015년은 강화학습 광풍이 불던 시기입니다. 구글 딥마인드가 아타리 게임에서 DQN으로 사람을 넘어섰고, AlphaGo가 바둑에서 돌파를 이루면서 강화학습에 엄청난 기대가 쏠렸어요. 그 상황에서 르쿤은 강화학습만으로는 인간 수준에 절대 못 간다고, 자기지도학습이 핵심이 되어야 한다고 주장한 거죠.

[말줄임] 결과적으로? 르쿤의 예측이 맞았어요. LLM의 성공 공식이 정확히 이 케이크 구조를 따릅니다. 자기지도 사전학습(다음 토큰 예측) $\to$ 지도학습 fine-tune $\to$ RLHF 강화학습 정렬. 케이크 그대로죠.

그런데 아이러니하게도, 이 성공이 \alert{비전이 아닌 언어에서 먼저 일어났습니다}. 왜 그랬을까요? 다음 파트에서 그 이유를 보고, 그다음에 영상에서 왜 어려웠는지를 봅시다.
\end{script}
\end{frame}

% ============================================================
\section{Part 3. LLM의 탄생: 다음 토큰 예측}

\begin{frame}{[8] 트랜스포머와 GPT-1 (2018)}
\begin{script}
OpenAI는 2015년에 설립됐을 때 처음에는 강화학습에 집중했어요. 그런데 2017년에 구글에서 나온 새로운 신경망 구조, \textbf{트랜스포머}가 등장하면서 흐름이 바뀝니다.

트랜스포머는 원래 기계 번역을 위해 설계됐어요. 한 언어의 텍스트 덩어리를 다른 언어의 텍스트 덩어리로 옮기는 거죠. 그런데 OpenAI의 알렉 래드퍼드가 여기서 흥미로운 변형을 시도했어요. 번역 대신, 훨씬 단순한 자기지도 방식으로 바꿔본 겁니다.

아이디어는 이렇습니다. 텍스트 시퀀스를 받아서, 마지막 토큰 하나만 빼고 나머지를 입력으로 주고, 그 마지막 토큰이 뭔지를 맞히게 합니다. 이걸 책 7,000권 규모의 데이터로 학습시킨 게 GPT-1이에요.

[말줄임] 결과가 놀라웠어요. 9개의 언어 벤치마크에서 새 SOTA를 갈아치웠는데, 그중에는 고등학교 수준의 독해 문제도 있었어요. 각 과제에 맞춰 따로 설계·학습된 전문 모델들을 모두 누른 거죠. GPT-1이 뭔가 더 일반적인 언어 이해를 학습했다는 증거였습니다. 이게 LLM 시대의 진정한 출발점이었어요.
\end{script}
\end{frame}

% ============================================================
\begin{frame}{[9] GPT 스케일링: 2012$\to$2020}
\begin{script}
GPT-1이 열어준 가장 중요한 문은 \textbf{스케일링}입니다. 사람이 만든 라벨이 필요 없으니, 인터넷에 있는 모든 텍스트를 학습 데이터로 쓸 수 있게 됐어요.

수치로 보면 극적입니다. 2012년 AlexNet은 약 100만 개의 예제로 학습됐어요. 2020년 GPT-3는 수천억 개의 토큰으로 학습됐습니다. 파라미터 수도 AlexNet의 6천만 개에서 GPT-3의 1750억 개로 수천 배 늘었어요.

[말줄임] 흥미로운 점은, 이 학습 패러다임이 르쿤이 케이크 슬라이드에서 그린 것과 정확히 일치한다는 거예요. 방대한 자기지도 사전학습 $\to$ 지도 fine-tune $\to$ 강화학습 정렬. 르쿤이 예언한 케이크를 LLM 팀이 언어에서 먼저 구워낸 거죠.

그런데 여기서 중요한 걸 짚어야 해요. 다음 토큰 예측 자체는 \textbf{프록시 태스크}예요. 겉으로는 자동완성기를 학습시키는 것처럼 보이지만, 실제로 우리가 얻는 건 모델 내부의 \textbf{표현}입니다. 언어 구조, 인과관계, 사실 지식, 논리 추론이 그 표현 안에 녹아있어요. 그 표현이 있으니까 fine-tune 몇 번으로 강력한 AI 비서가 되는 거죠.

이 통찰이 JEPA로 이어집니다. 꼭 생성 모델이어야 좋은 표현을 얻는 게 아니라면, 생성을 건너뛰고 표현만 학습할 수 없을까?
\end{script}
\end{frame}

% ============================================================
\begin{frame}{[10] 사전학습이 배운 것의 본질}
\begin{script}
이 슬라이드에서 핵심 통찰 하나를 짚고 갑시다. LLM 성공의 진짜 비밀이 뭐냐는 겁니다.

표면적으로는 ``다음 단어가 뭐지?''라는 단순한 자동완성 과제예요. 그런데 이걸 정말 잘 풀려면 모델이 내부적으로 무엇을 알아야 할까요? 인과관계, 사실 지식, 논리 추론, 코딩 능력, 수학 추론, 심지어 유추까지요. 이게 다 내부 표현으로 녹아들어 있어요.

[말줄임] 그러면 이런 질문을 할 수 있어요. \textit{``꼭 텍스트를 생성해야 이 표현을 얻는 건가? 생성 없이도 그에 버금가는 표현을 얻을 수 없을까?''}

이게 JEPA가 던지는 근본 질문이에요. 그리고 이 질문을 먼저 비전(이미지, 영상) 도메인에서 풀려고 했을 때 어떤 어려움을 만났는지가 다음 파트의 이야기입니다.

LLM에서 자기지도 생성이 그렇게 잘 된 건 언어의 특성 덕분인데, 영상은 그 특성이 근본적으로 달라요.
\end{script}
\end{frame}

% ============================================================
\section{Part 4. 영상의 흐릿함 문제}

\begin{frame}{[11] 픽셀 생성: 왜 영상은 텍스트와 다른가}
\begin{script}
LLM이 텍스트에서 성공하기 한참 전부터, 르쿤을 포함한 연구자들은 같은 자기지도 생성 방식을 영상에 적용해 봤어요.

가장 단순한 구현은 이렇습니다. 신경망에 비디오 프레임 시퀀스를 입력으로 주고, 다음 프레임의 픽셀 값을 예측하게 합니다. GPT가 다음 토큰을 예측하듯이요.

문제가 뭐냐면, 언어와 영상의 \textbf{출력 공간이 근본적으로 다르다}는 거예요. 언어 모델의 출력은 고정된 어휘에서 하나를 고르는 거예요. GPT-2 기준으로 50,257개의 이산 토큰 중 하나죠. 이 토큰들에 각각 확률을 부여할 수 있어요.

[말줄임] 영상에서는요? Full HD 영상 한 프레임의 픽셀 수는 1920 곱하기 1080 곱하기 3, 각 픽셀이 0에서 255 사이 값을 가집니다. 가능한 다음 프레임의 수는 대략 $10^{15,000,000}$이에요. 관측 가능한 우주의 원자 수인 $10^{80}$을 우습게 넘기죠. 이걸 토큰화해서 각각 노드를 두는 건 아예 불가능합니다.

그러면 어떻게 하느냐? 픽셀 강도를 실수값으로 직접 예측하도록 합니다. 거기서 진짜 문제가 생겨요.
\end{script}
\end{frame}

% ============================================================
\begin{frame}{[12] 평균화 = 흐릿함}
\begin{script}
픽셀 강도를 실수값으로 예측할 때의 핵심 문제는 \textbf{불확실성을 어떻게 다루는가}입니다.

비유를 들어볼게요. 공이 화면 왼쪽으로 튀는 예시와 오른쪽으로 튀는 예시가 학습 데이터에 섞여 있다고 합시다. LLM이 같은 상황에 처하면 어떻게 할까요? ``왼쪽'' 토큰의 확률을 높이고, 동시에 ``오른쪽'' 토큰의 확률을 독립적으로 높일 수 있어요. 두 가능성이 분리된 토큰으로 존재하니까요.

영상 모델은요? 단 하나의 픽셀 배열을 직접 출력해야 합니다. 손실을 최소화하는 가장 쉬운 답은 두 가능한 결과의 픽셀값들을 평균 내는 거예요. 그러면 어떻게 되냐? 공이 왼쪽에 반쯤, 오른쪽에 반쯤 걸쳐 있는 흐릿하고 비현실적인 이미지가 됩니다.

[말줄임] 이걸 그냥 ``구현 문제''로 치부할 수도 있는데, 이게 사실 더 깊은 질문을 던져요. \textit{꼭 모델이 생성형이어야 하는가?} GPT에서 진짜로 중요한 건 텍스트를 생성하는 능력 자체가 아니라, 그 학습 과정에서 얻은 \textbf{내부 표현}이잖아요. 그렇다면 생성 없이도 좋은 표현을 얻을 수 있지 않을까요? 그게 결합 임베딩 아이디어예요.
\end{script}
\end{frame}

% ============================================================
\section{Part 5. 결합 임베딩과 샴 신경망}

\begin{frame}{[13] 비생성적 표현 학습: 아이디어의 전환}
\begin{script}
르쿤이 2017--18년쯤부터 깨닫기 시작한 게 있어요. 이미지 표현을 가장 잘 배우는 시스템은 \textbf{생성하지 않는 시스템}이라는 거죠.

아이디어는 이렇습니다. 같은 장면을 찍은 두 이미지를 가져오거나, 한 이미지를 어떤 식으로 변형해서 또 하나를 만들어요. 둘 다 인코더에 통과시키고, 두 표현이 의미적으로 같은 것을 가리키니 최대한 비슷해지도록 시스템을 학습시킵니다.

[말줄임] 그런데 이건 사실 르쿤이 1990년대에 이미 한 적 있는 발상이에요. 슬라이드에서 언급한 \textbf{샴 신경망(Siamese network)}입니다.

슬라이드를 보면서 설명할게요. 두 개의 같은 가중치를 공유하는 네트워크, 각각 이미지 하나씩을 입력으로 받아서 임베딩 벡터를 출력해요. 이미지를 생성하는 게 아니에요. 그냥 숫자 벡터 하나씩. 그 두 벡터의 거리를 재서 학습합니다.

이 방식의 핵심은, 흐릿함 문제를 원천 회피한다는 거예요. 픽셀 공간에서 예측하지 않고 추상 표현 공간에서 거리를 재니까요. 다음 슬라이드에서 샴 신경망의 원조 이야기를 봅시다.
\end{script}
\end{frame}

% ============================================================
\begin{frame}{[14] 샴 신경망: 1990년대의 원조}
\begin{script}
샴 신경망이 탄생한 맥락이 흥미롭습니다. 1990년대 초, 르쿤과 Bell Labs 동료들이 \textbf{위조 서명을 잡아내려고} 만든 구조예요.

작동 방식을 봅시다. 두 장의 서명 이미지를 가중치를 공유하는 두 개의 복제 네트워크에 각각 통과시킵니다. 네트워크는 서명 이미지를 생성하지 않아요. 대신 \textbf{임베딩 벡터}라고 부르는 숫자 벡터를 출력합니다.

학습은 두 종류의 예시로 해요. 양성 쌍: 기준 서명과 같은 사람의 진짜 서명 쌍. 음성 쌍: 기준 서명과 위조 서명 쌍. 양성 쌍에 대해서는 두 임베딩이 비슷해지도록, 음성 쌍에 대해서는 달라지도록 학습합니다.

[말줄임] 이 과정에서 서명의 무엇이 ``진짜냐 위조냐''를 결정하는 특성들이 임베딩 벡터 안에 자연스럽게 인코딩됩니다. 서명 이미지를 다시 만들지 않고도요. 이게 결합 임베딩의 핵심 통찰이에요.

그렇다면 이미지와 영상에도 같은 방식을 적용할 수 있지 않을까? 같은 장면의 다른 뷰 두 개를 인코더에 통과시켜 비슷한 임베딩으로 매핑하도록 학습하면, GPT가 사전학습으로 유용한 표현을 얻듯이 일반적인 시각 표현을 얻을 수 있지 않을까?

문제는, 그렇게 단순할 리가 없다는 거였어요.
\end{script}
\end{frame}

% ============================================================
\section{Part 6. 표현 붕괴와 Barlow Twins}

\begin{frame}{[15] 표현 붕괴: 가장 게으른 답}
\begin{script}
샴 신경망 방식을 이미지에 적용하면 치명적인 함정이 있어요. 같은 이미지의 두 변형을 인코더에 넣고 임베딩이 비슷해지도록 학습시킨다고 했죠? 네트워크 입장에서 가장 손쉬운 답이 뭘까요?

\textbf{모든 입력에 대해 동일한 벡터를 출력하는 것}입니다. 예를 들어 어떤 이미지가 들어오든 전부 [1, 1, 1, ..., 1]을 출력하면, 같은 이미지의 두 변형은 당연히 두 벡터가 완전히 같으니까 유사도가 최대가 돼요. 손실 함수는 완전히 만족하지만, 모델이 배운 건 \alert{아무것도 없어요}. 이게 \textbf{표현 붕괴(representation collapse)}예요.

[말줄임] 처음에는 음성 쌍(negative pair)을 추가한 \textbf{대조 학습(contrastive learning)}으로 이를 피했어요. 다른 이미지들끼리는 임베딩이 달라지도록 학습시키는 거죠. 이건 실제로 잘 작동했어요. 문제는 스케일을 키울수록 골치 아파진다는 거였어요. 의미 있는 표현을 얻으려면 엄청난 계산량과 많은 음성 샘플이 필요하고, 최악의 경우 필요한 음성 샘플 수가 표현 차원에 대해 지수적으로 증가할 수 있거든요.

더 깔끔한 해법이 필요했고, 그 깨달음이 Barlow Twins에서 왔습니다.
\end{script}
\end{frame}

% ============================================================
\begin{frame}{[16] Barlow Twins: 뉴런 중복성을 줄여라}
\begin{script}
깨달음의 출처가 흥미롭습니다. Meta의 박사후연구원이었던 스테판 데니가 1961년의 이론신경과학자 \textbf{호러스 발로(Horace Barlow)}의 작업에서 힌트를 얻어왔어요.

발로의 가설은 이렇습니다. 동물과 사람의 시각 시스템에서 뉴런들은 서로 간의 \textbf{중복된 정보를 줄이는} 방향으로 작동한다. 이 아이디어를 신경망 출력에 적용해 보자는 게 데니의 제안이었어요.

슬라이드의 교차상관 행렬을 봅시다. 배치의 이미지들을 두 인코더에 넣으면, 각 인코더의 각 뉴런이 배치에 걸쳐서 활성값 시퀀스를 만들어요. 인코더 A의 $i$번 뉴런과 인코더 B의 $j$번 뉴런 사이의 피어슨 상관계수를 $(i,j)$ 성분으로 하는 행렬이 교차상관 행렬입니다.

[말줄임] 이 행렬이 \textbf{단위 행렬}이 되도록 학습하면 두 가지가 동시에 달성됩니다. 대각 성분이 1에 가까워지면 같은 이미지의 변형 쌍에 대해 대응하는 뉴런들이 일관되게 반응하는 거예요. 비대각 성분이 0에 가까워지면 서로 다른 뉴런들이 중복되지 않는 독립적인 정보를 인코딩하게 돼요. 그리고 이 비대각 조건이 표현 붕괴를 막아줍니다. 모든 뉴런이 같은 걸 출력하면 비대각 성분이 1이 되어서 손실이 커지거든요.
\end{script}
\end{frame}

% ============================================================
\begin{frame}{[17] Barlow Twins 손실 함수}
\begin{script}
손실 함수를 좀 더 구체적으로 봅시다.

두 항으로 구성됩니다. \textbf{불변성 항}은 대각 원소를 1로 당겨서, 같은 이미지의 변형 쌍에 대해 같은 표현을 만들게 강제합니다. 이게 없으면 네트워크가 원본과 변형에 전혀 다른 표현을 할당해버려요. \textbf{중복성 항}은 비대각 원소를 0으로 당겨서, 뉴런들이 서로 다른 정보를 인코딩하도록 강제합니다. 이게 바로 발로의 아이디어예요.

[말줄임] 이 방식의 핵심 장점은 \textbf{음성 샘플이 필요 없다}는 거예요. 대조 학습은 음성 쌍이 있어야 했잖아요. 서로 다른 이미지들을 대량으로 모아서 ``이건 달라''라고 알려줘야 했죠. Barlow Twins는 그게 없어도 됩니다. 손실 함수 안의 중복성 항이 그 역할을 수학적으로 대체하거든요.

실제로 결과도 좋았어요. 표현 붕괴를 피하면서 동시에 강력한 내부 표현을 학습하는 데 성공했습니다. 다음 슬라이드에서 그 표현의 품질을 어떻게 측정하고, 실제로 얼마나 좋은지를 봅시다.
\end{script}
\end{frame}

% ============================================================
\section{Part 7. 표현 품질 측정과 DINO}

\begin{frame}{[18] 선형 프로브: 표현 품질을 재는 도구}
\begin{script}
자기지도로 학습한 표현이 정말 쓸 만한지를 어떻게 검증할까요? 가장 표준적인 방법이 \textbf{선형 프로브}입니다.

방법은 이렇습니다. Barlow Twins 같은 자기지도 방법으로 학습한 인코더를 가져옵니다. 이 인코더 위에 딱 \textbf{한 층의 선형 분류기}만 추가해요. 그리고 그 한 층만 ImageNet 라벨로 지도학습시킵니다. 인코더 본체는 \textbf{동결(frozen)}해서 절대 건드리지 않아요.

이게 측정하는 건 순수하게 ``동결된 표현만으로, 클래스들이 선형적으로 분리 가능한 구조가 잡혀 있는가?''입니다. 표현이 나쁘면 선형 분류기 하나로 아무리 해봤자 정확도가 안 나와요.

[말줄임] 결과를 보면 인상적이에요. 슬라이드의 표를 같이 봅시다. 원조 AlexNet이 완전 지도학습으로 59.3%를 냈어요. 2012년에요. Barlow Twins는 \textbf{라벨 한 개도 보지 않고} 자기지도학습만으로 73.2%를 달성했어요. 거기다가 선형 프로브 한 층만 추가했는데요.

다만 2020년에 구글 팀이 완전 지도학습으로 88.6\%를 달성한 상태라 아직 격차가 있었어요. 그 격차를 메워가는 게 DINO 계열입니다.
\end{script}
\end{frame}

% ============================================================
\begin{frame}{[19] VICReg $\to$ DINO: 자기지도학습이 따라붙다}
\begin{script}
Barlow Twins 이후 이 방향이 옳다는 게 분명해지면서, 여러 팀이 비슷한 계열의 방법을 발전시켰어요.

Meta에서 나온 \textbf{VICReg}는 Barlow Twins의 단순화 버전이에요. 비슷한 원리를 좀 더 직관적인 형태로 정리한 거죠. 동시에 FAIR Paris의 동료들이 연구하고 있었는데, 그게 \textbf{DINO}로 알려지게 됩니다.

DINO는 지식 증류(knowledge distillation)를 이용한 결합 임베딩 방법이에요. v1, v2를 거쳐 2025년 8월에 나온 DINO v3에서 결정적인 이정표를 세웠습니다. ImageNet top-1 정확도 \alert{88.4\%}. 지도학습 SOTA인 88.6\%와 사실상 같은 수준이에요.

[말줄임] 저자들이 논문에서 직접 이렇게 말했습니다. \textit{``자기지도 모델이 이미지 분류에서 지도학습 모델과 견줄 만한 결과를 낸 것은 이번이 처음이다.''} 굉장히 중요한 선언이에요.

더 인상적인 건 질적인 특성입니다. DINO는 이미지 패치마다 임베딩 벡터를 출력하는데, 어떤 사람의 손 위치 패치의 임베딩을 뽑아서 다른 패치들과의 유사도를 색으로 표시하면, 라벨 없이도 손과 배경을 깔끔하게 분리해냅니다. 별도의 세그멘테이션 학습 없이요. 이게 지금 우리가 다루고 있는 표현 학습의 힘이에요.
\end{script}
\end{frame}

% ============================================================
\section{Part 8. JEPA 아키텍처: 세계 모델로}

\begin{frame}{[20] 왜 세계 모델인가: 운전하는 십대}
\begin{script}
자기지도 결합 임베딩이 비전에서 잘 작동한다는 게 증명됐습니다. 그런데 르쿤은 여기서 멈추지 않아요. Barlow Twins, VICReg, DINO v1의 성공을 등에 업고 2022년에 60쪽짜리 입장 논문을 내놓습니다. 제목: \textit{A Path Towards Autonomous Machine Intelligence}.

이 논문의 출발점이 바로 슬라이드에서 본 ``운전하는 십대'' 예시예요. 르쿤의 말로 들어봅시다. \textit{``수백만 시간의 주행 데이터로 우리가 만들어 낼 수 있는 건 테슬라의 레벨 2 시스템 정도예요. 레벨 3, 4, 5는 근처에도 못 갔습니다. 그런데 17살짜리는 몇 시간 연습으로 운전을 배워요.''}

이 차이가 어디서 오느냐. 르쿤의 답은 \textbf{세계 모델}입니다. 인간과 동물은 물리 세계가 어떻게 작동하는지에 대한 내부 모델을 갖고 있고, 그 모델을 이용해서 자기 행동의 결과를 \alert{사전에 시뮬레이션}하고 계획합니다. 17살짜리가 핸들을 돌리기 전에 ``이렇게 하면 차가 저쪽으로 갈 것''을 머릿속으로 예측할 수 있는 거예요.

[말줄임] 현재의 AI에는 이게 없다는 게 르쿤의 핵심 주장이에요. LLM도, VLA도, 자기 행동의 결과를 사전에 예측하는 능력이 없습니다. 행동을 출력하고 나서야 다음 상태를 관측할 수 있는 거죠. 르쿤은 이를 ``après moi, le déluge'', 즉 ``내 뒤에 대홍수가 오든 말든''이라고 표현했어요.
\end{script}
\end{frame}

% ============================================================
\begin{frame}{[21] JEPA 아키텍처 상세}
\begin{script}
이제 JEPA 구조를 구체적으로 봅시다.

슬라이드의 다이어그램을 같이 봐요. 기본 JEPA는 이렇게 생겼어요. 현재 관측 $x$를 인코더에 통과시켜 임베딩 $s_x$를 얻습니다. 목표나 미래 관측 $y$도 인코더에 통과시켜 $s_y$를 얻습니다. 그리고 예측기가 $s_x$에서 $s_y$를 예측하도록 학습합니다.

여기서 중요한 설계 요소가 \textbf{stop\_grad}입니다. 목표 인코더는 역전파로 직접 학습되지 않아요. 대신 EMA(지수이동평균)로 천천히 업데이트됩니다. 이게 왜 중요하냐면, 두 인코더를 동시에 역전파로 학습시키면 표현 붕괴가 다시 일어날 수 있거든요. stop\_grad가 이를 막는 핵심 설계예요.

[말줄임] 그리고 결정적인 확장이 있어요. 예측기에 \textbf{행동(action)을 조건으로 줄 수 있다}는 거예요. 시각 $t$의 임베딩 $s_t$와 그 시점에 취한 행동 $a_t$를 예측기에 넣으면, 예측기가 시각 $t+1$의 임베딩 $s_{t+1}$을 예측하도록 학습됩니다. 이 순간 JEPA는 \textbf{세계 모델}이 돼요. ``이 행동을 하면 세계 상태가 어떻게 변하는가''를 임베딩 공간에서 예측하는 모델이 되는 거죠.
\end{script}
\end{frame}

% ============================================================
\begin{frame}{[22] JEPA가 세계 모델이 되는 순간: 행동 조건화}
\begin{script}
이걸 로봇에 적용하면 어떻게 되는지를 V-JEPA 2를 통해 봅시다.

팀은 로봇 팔과 환경을 찍은 비디오 시퀀스를 가져와요. 각 시점의 프레임을 인코더에 통과시켜 임베딩을 얻습니다. 그리고 예측기가 현재 임베딩과 그 시점에 로봇 팔에 보낸 \textbf{제어 신호}를 받아서, 다음 프레임의 임베딩을 예측하도록 학습합니다.

학습이 되면 예측기는 ``어떤 제어 신호가 임베딩 공간에서 로봇의 상태를 어떻게 바꾸는지''를 알게 됩니다. 이게 바로 르쿤이 말하는 세계 모델이에요.

[말줄임] 이 세계 모델을 계획에 쓰는 방식을 봅시다. 목표 상태, 예를 들어 ``컵이 저기 있는'' 장면의 이미지를 인코더에 통과시켜 목표 임베딩을 얻습니다. 그리고 제어 알고리즘이 다양한 가상의 행동 시퀀스를 세계 모델에 넣어보면서, 예측되는 미래 상태가 목표 임베딩에 가장 가까워지는 행동 시퀀스를 찾아냅니다.

이게 \textbf{고전적 최적 제어}예요. 르쿤이 직접 ``1950년대 소련, 1960년대 서구로 거슬러 올라가는 아이디어''라고 말했어요. 새로운 건 그 위에 머신러닝과 학습된 추상 표현이 더해졌다는 거죠.
\end{script}
\end{frame}

% ============================================================
\begin{frame}{[23] 세계 모델 = 고전 최적 제어 + 학습된 표현}
\begin{script}
이 슬라이드에서 VLA와 JEPA를 직접 비교해 봅시다.

VLA, 즉 $\pi_0$ 같은 방식에서는 이미지와 텍스트를 받아서 행동을 \textbf{직접 출력}합니다. 내부에서 어떤 연산을 하든, 출력은 행동 토큰이나 관절 궤적이에요.

JEPA 세계 모델에서는 다릅니다. 가능한 행동 시퀀스들을 세계 모델에 넣어보면서, ``이 행동을 하면 세계가 어떻게 될까''를 먼저 예측하고, 그 예측 결과 중 목표에 가장 가까운 것을 선택하는 \textbf{최적화(search)} 과정이 있습니다.

[말줄임] 르쿤의 표현이 인상적이에요. VLA와 LLM에 대해 이렇게 말했습니다. \textit{``LLM도 VLA도 자기 행동의 결과를 사전에 예측하지 못한다. 그냥 행동부터 하고 나서 내 뒤에 대홍수가 오든 말든인 거다.''} 반면 JEPA 세계 모델은 \textit{``행동하기 전에 그 결과를 시뮬레이션하고 계획한다''}.

표에서 보이듯이 두 접근이 근본적으로 다른 가정 위에 서 있어요. 어느 쪽이 옳은지는 아직 열려 있는 질문입니다. 하지만 이 두 가정을 명확히 이해하는 게 앞으로의 연구 방향을 잡는 데 중요해요.
\end{script}
\end{frame}

% ============================================================
\section{Part 9. JEPA가 LLM을 대체할 것인가}

\begin{frame}{[24] LeCun 인터뷰: 열린 대화}
\begin{script}
르쿤의 입장을 인터뷰 형식으로 정리한 슬라이드를 봅시다.

진행자가 묻습니다. ``JEPA가 언젠가 LLM을 대체할 거라고 보세요?'' 르쿤의 첫 답변은 의외로 신중해요. ``당분간은 서로 다른 문제를 풀 겁니다.'' 진행자가 다시 확인하자, 그제야 ``네, 결국은 대체합니다. LLM은 언어에는 능숙하지만 사실 그게 전부거든요''라고 말합니다.

[말줄임] 그의 핵심 주장을 다시 정리하면 이렇습니다. 신뢰할 수 있는 에이전트 시스템이 되려면 세 가지가 필요해요. 첫째, 행동 결과를 사전에 예측할 수 있어야 한다. 둘째, 그 예측을 바탕으로 행동 시퀀스를 계획할 수 있어야 한다. 셋째, 이 과정에서 안전 가드레일을 보장할 수 있어야 한다.

진행자가 이를 받아서 이렇게 정리해요. ``그러면 추론 과정 자체가 단순한 자기회귀 예측이 아니라 탐색(search)이 되는 거네요.'' 르쿤: ``맞아요. 그게 바로 세계 모델입니다.''

이게 오늘 강의의 핵심 메시지예요. LLM은 매우 강력하지만, 세계를 모델링하는 능력이 없다는 근본적 한계가 있고, JEPA는 그 한계를 우회하는 다른 길을 제시한다는 거죠.
\end{script}
\end{frame}

% ============================================================
\begin{frame}{[25] 토론거리}
\begin{script}
오늘도 마무리는 토론으로 합시다. 다섯 가지 질문을 드릴게요.

\textbf{첫 번째}는 Barlow Twins의 생물학적 타당성에 관한 거예요. 발로의 가설이 1961년에 나왔는데, 현재 신경과학에서 이게 얼마나 지지받고 있는지, 그리고 Sparse coding이나 ICA 같은 다른 신경과학 기반 표현 학습 이론과 어떻게 연결되는지 생각해 보세요.

\textbf{두 번째}는 stop\_grad에 관한 건데요. stop\_grad 없이 JEPA를 학습시키면 왜 붕괴하는지, 이게 심층강화학습에서 나오는 \textbf{Bootstrapping 문제}와 어떻게 연결되는지 생각해 봐요.

[말줄임] \textbf{세 번째}: DINO v3가 ImageNet 88.4\%를 달성했지만, 어떤 태스크에서 자기지도학습이 아직 지도학습보다 열위인지 조사해 보세요. 왜 그 격차가 생기는지.

\textbf{네 번째}: VLA와 JEPA World Model을 하이브리드로 결합하는 방법을 설계해 보세요. 예를 들어 $\pi_0$의 Action Expert 대신 JEPA predictor를 넣으면 어떻게 될까요?

\textbf{다섯 번째}는 APRL 연구와 직결되는 질문이에요. SLAM은 이미 세계 모델의 일종으로 볼 수 있어요. 환경 지도를 유지하고 업데이트하는 거잖아요. JEPA 스타일 세계 모델과 classical SLAM의 차이가 뭐고, 어떻게 보완할 수 있을지 생각해 봐요.
\end{script}
\end{frame}

% ============================================================
\begin{frame}{[26] 핵심 정리}
\begin{script}
오늘 긴 여정을 했습니다. 핵심 다섯 가지로 정리할게요.

\textbf{첫째, LeCun의 도발}: LLM과 VLA는 세계 모델이 없어서 자기 행동의 결과를 사전에 예측하지 못한다. 이게 신뢰할 수 있는 에이전트의 근본 한계다.

\textbf{둘째, 역사의 흐름}: CNN의 라벨 의존성 문제 $\to$ 자기지도학습 필요성 인식 $\to$ LLM에서 언어로 먼저 폭발(다음 토큰 예측이 이산 어휘 덕분에 잘 작동). 영상은 연속 픽셀 공간 때문에 같은 방법이 흐릿함 문제를 일으킴.

\textbf{셋째, 결합 임베딩의 발전}: 샴 신경망(1990s) $\to$ 표현 붕괴 문제 $\to$ 대조 학습(한계) $\to$ Barlow Twins(교차상관 행렬 = 단위 행렬) $\to$ DINO v3(자기지도 = 지도학습 수준).

\textbf{넷째, JEPA}: 인코더 두 개 + 예측기로 latent space에서 예측. 행동 조건화로 세계 모델로 발전. 고전 최적 제어 + 학습된 표현.

\textbf{다섯째, 미래}: 단기는 공존, 장기는 대체라는 게 르쿤의 입장. 세계 모델 없는 AI에 르쿤이 내기를 걸지 않는 이유는 명확합니다. 행동 전에 결과를 예측할 수 없으면, 진정한 의미의 계획도, 안전도 보장할 수 없다는 거죠.
\end{script}
\end{frame}

% ============================================================
\begin{frame}{[27] 참고 문헌}
\begin{script}
오늘 강의에서 직접 다룬 핵심 논문들을 정리합니다.

Barlow Twins(2021), VICReg(2022), DINO 계열은 결합 임베딩 자기지도학습의 핵심 페이퍼예요. 특히 Barlow Twins 논문은 손실 함수 유도 과정이 깔끔하게 서술되어 있어서 읽기 좋습니다.

I-JEPA(2023)와 V-JEPA(2024), V-JEPA 2(2025)는 JEPA 아키텍처를 이미지, 비디오, 로봇 제어에 적용한 논문들이에요. 다음 주에 V-JEPA 2 구현을 상세히 다룰 예정이니 미리 읽어오면 좋습니다.

르쿤의 2022년 입장 논문은 오늘 강의의 철학적 토대가 됐어요. 길지만 읽어볼 가치가 있습니다. 특히 세계 모델, 계획, 추론에 대한 섹션을 중심으로 보세요.

[말줄임] 다음 주에는 V-JEPA 2의 구현 상세를 들여다보고, VLA 방법들과 실제 성능을 비교해볼 거예요. 특히 로봇 제어 벤치마크에서 두 패러다임이 어떻게 맞부딪히는지를 봅시다. 오늘 수고하셨습니다.
\end{script}
\end{frame}

\end{document}
